% ══════════════════════════════════════════════════════════════════════════════
%  BIODIVERSITY AND ECOSYSTEMS — A5 LESSON BOOKLET, LESSONS 9–11
% ══════════════════════════════════════════════════════════════════════════════

% ══════════════════════════════════════════════════════════════════════════════
%  LESSON 9 — INTRODUCED SPECIES
% ══════════════════════════════════════════════════════════════════════════════
\setcounter{section}{8}
\sectiondiv{Introduced Species}{Native, introduced and invasive \textbullet{} Four ways they cause harm \textbullet{} Impact chains}\label{sec:9}
\lessoncodes{VC2S8U04 \textbullet{} VC2S8H03}

\vocabhead
\begin{ctab}
\begin{tabular}{|SW{3.4cm}|W{7.0cm}|}
  \hline\hd & \thead{Word} & \thead{Definition (copy from the slide)} \\ \hline
  & \scaf{\textbf{Native species}} & \scaf{A species that has always lived in an area, without being brought by people.} \\ \hline
  & \textbf{Introduced species} & \\ \hline
  & \textbf{Invasive species} & \\ \hline
  & \textbf{Predation} & \\ \hline
  & \textbf{Competition} & \\ \hline
  & \textbf{Biological control} & \\ \hline
\end{tabular}
\end{ctab}

% ─── PART A ───────────────────────────────────────────────────────────────────
\lessonpart{A}{Reading and comprehension}{Teal}\label{read:9}

\begin{readingbox}{Reading --- Twenty-Four Rabbits}
In 1859 a farmer near Geelong released 24 rabbits so he could hunt them. With plenty of grass and few predators, they bred so fast that within 50 years there were hundreds of millions across Australia. They stripped the land of plants and left the soil bare.

Rabbits are an \textbf{introduced species}: brought to a new place by people, on purpose or by accident. A \textbf{native species}, like the eastern grey kangaroo, has always lived here. Not every introduced species causes trouble, but one that spreads quickly and does real harm is called an \textbf{invasive species}.

Invasive species do harm in four main ways. Through \textbf{predation} they hunt native animals that never evolved to escape them: foxes and cats kill bandicoots and ground-nesting birds. Through \textbf{competition} they take food, water or shelter that native species need: rabbits eat the grass kangaroos rely on. They change habitats, as blackberry does when it smothers creek banks, and they spread diseases.

People fight back with fences, baiting, trapping and \textbf{biological control}: using a living thing, such as a virus, to control a pest. In 1950 the myxoma virus was released and killed most of Australia's rabbits, though survivors slowly became resistant.
\end{readingbox}

\rcreset
\rcq{Why did rabbits spread so quickly in Australia?}{2.5cm}
\rcq{What is the difference between an introduced species and an invasive species?}{2.5cm}
\rcq{Name the four ways invasive species cause harm.}{2.5cm}
\rcq{What is biological control? Give an example.}{2.5cm}
\rcq{Why did the myxoma virus not wipe out rabbits for good?}{2cm}

% ─── PART B ───────────────────────────────────────────────────────────────────
\lessonpart{B}{Theory and examples}{Blue}\label{theory:9}

\subsection{Four Ways of Causing Harm}\label{sub:9.1}

\begin{ctab}
\begin{tabular}{|L{2.3cm}|L{3.8cm}|L{3.9cm}|}
  \hline\thead{Harm} & \thead{What happens} & \thead{Victorian example} \\ \hline
  Predation & Eats native animals with no defences against it & Red fox killing bandicoots \\ \hline
  Competition & Takes food, water or shelter natives need & Rabbits eating kangaroos' grass \\ \hline
  Habitat change & Alters the environment & Blackberry smothering creek banks \\ \hline
  Disease & Carries germs natives cannot fight & Cats spreading toxoplasmosis \\ \hline
\end{tabular}
\end{ctab}

\subsection{Impact Chains}\label{sub:9.2}

\begin{workedex}{European carp}
Carp were released into Victorian rivers in the 1960s. They feed by sucking up mud from the riverbed.

\smallskip
\textbf{Step 1:} feeding stirs up the mud, making the water cloudy and releasing nutrients.\\
\textbf{Step 2:} less light reaches the riverbed, so water plants cannot photosynthesise and die.\\
\textbf{Step 3:} extra nutrients feed algae; when the algae die and rot, oxygen in the water drops.\\
\textbf{Step 4:} native fish such as Murray cod lose shelter, food and oxygen; platypus lose the insect larvae they eat.

\smallskip
One introduced species changed a whole river through a chain of effects: an \textbf{impact chain}.
\end{workedex}

\begin{wedo}{Red foxes}
Red foxes were brought to Victoria in the 1870s for hunting. Trace an impact chain of at least three steps, starting with the foxes eating bandicoots and ground-nesting birds.
\answerbox{5cm}
\end{wedo}

\begin{countrybox}{First Nations perspectives --- Loss of Country and culture}
Before 1835, bettongs, quolls, bandicoots and bush stone-curlews were common on Wurundjeri Woi-wurrung Country. Within decades, sheep, cattle, rabbits, foxes, cats and blackberry arrived together and devastated them. This was not only an environmental loss: many of these animals are totems for particular families, so their loss is a loss of identity and culture. Restoration projects led with Wurundjeri Elders are cultural restoration too.
\end{countrybox}

% ─── PART C ───────────────────────────────────────────────────────────────────
\lessonpart{C}{Practice}{Amber}\label{prac:9}

\task Sort each species as native or introduced to Victoria. For introduced species, give one way it could cause harm: eastern grey kangaroo, European rabbit, red fox, tawny frogmouth, blackberry, common starling, platypus, European carp.
\begin{ctab}
\begin{tabular}{|SW{3.0cm}|W{2.2cm}|W{4.8cm}|}
  \hline\hd & \thead{Species} & \thead{Native or introduced} & \thead{How it could cause harm} \\ \hline
  & \scaf{European rabbit} & \scaf{Introduced} & \scaf{Competition: eats grass native grazers need} \\ \hline
  & & & \\ \hline
  & & & \\ \hline
  & & & \\ \hline
  & & & \\ \hline
  & & & \\ \hline
  & & & \\ \hline
  & & & \\ \hline
\end{tabular}
\end{ctab}

\needspace{10\baselineskip}
\task Native mammals such as bandicoots are easily killed by foxes, but rabbits, also introduced, survive alongside foxes. Suggest why.
\answerbox{5cm}

\needspace{10\baselineskip}
\task Cats were brought to Australia as pets. Explain how a pet can become an invasive species.
\answerbox{5cm}

% ─── PART D ───────────────────────────────────────────────────────────────────
\lessonpart{D}{Above level}{Violet}\label{ext:9}

\begin{yr8box}
These questions go beyond the Year 7 standard.

\medskip
\begin{keepext}
\exttask Evaluate the myxoma virus as a biological control for rabbits: describe how it worked, the evidence that it worked, and one risk of using biological control. (Cane toads are one example of a biological control that went wrong.)
\answerbox{6cm}
\end{keepext}
\begin{keepext}
\exttask Animals on islands are especially at risk from introduced predators. Explain why, and give a Victorian example.
\answerbox{5cm}
\end{keepext}
\end{yr8box}

% ══════════════════════════════════════════════════════════════════════════════
%  LESSON 10 — HUMAN IMPACTS ON ECOSYSTEMS
% ══════════════════════════════════════════════════════════════════════════════
\setcounter{section}{9}
\sectiondiv{Human Impacts on Ecosystems}{Habitat loss \textbullet{} Pollution and eutrophication \textbullet{} Climate change \textbullet{} Threats that combine}\label{sec:10}
\lessoncodes{VC2S8U04 \textbullet{} VC2S8H03}

\vocabhead
\begin{ctab}
\begin{tabular}{|SW{3.4cm}|W{7.0cm}|}
  \hline\hd & \thead{Word} & \thead{Definition (copy from the slide)} \\ \hline
  & \scaf{\textbf{Habitat loss}} & \scaf{When the place a species lives is destroyed or changed so it can no longer live there.} \\ \hline
  & \textbf{Fragmentation} & \\ \hline
  & \textbf{Pollution} & \\ \hline
  & \textbf{Eutrophication} & \\ \hline
  & \textbf{Over-harvesting} & \\ \hline
  & \textbf{Endangered} & \\ \hline
\end{tabular}
\end{ctab}

% ─── PART A ───────────────────────────────────────────────────────────────────
\lessonpart{A}{Reading and comprehension}{Teal}\label{read:10}

\begin{readingbox}{Reading --- Two-Thirds Gone}
Since European settlement, Victoria has cleared about two-thirds of its native vegetation for farms, towns and roads. The biggest threat to wildlife is \textbf{habitat loss}. When a forest is cleared, the animals that lived there have nowhere to go.

Often the bush is not all cleared, but cut into pieces by roads and paddocks. This \textbf{fragmentation} leaves small, separate patches. Animals in one patch cannot reach mates or food in another, and the edges of each patch get hotter, windier and weedier.

\textbf{Pollution} adds harmful substances to the environment: plastic, pesticides, and fertiliser washed into rivers. Too much fertiliser in water causes \textbf{eutrophication}: algae grow wildly, then die and rot, and the bacteria that rot them use up the oxygen, killing fish.

Climate change brings hotter summers, less snow and bigger fires. \textbf{Over-harvesting} takes fish, animals or plants faster than they can breed. These threats rarely act alone. A species facing several at once can become \textbf{endangered}: at serious risk of dying out.
\end{readingbox}

\rcreset
\rcq{What is the biggest threat to wildlife in Victoria?}{2cm}
\rcq{What is fragmentation, and why is it harmful?}{2.5cm}
\rcq{Explain how fertiliser can kill fish in a lake.}{3cm}
\rcq{Name two effects of climate change on Victorian ecosystems.}{2cm}
\rcq{What does it mean for a species to be endangered?}{2cm}

% ─── PART B ───────────────────────────────────────────────────────────────────
\lessonpart{B}{Theory and examples}{Blue}\label{theory:10}

\subsection{Fragmentation}\label{sub:10.1}

\begin{center}
\begin{tikzpicture}[font=\small]
  \fill[black!25, draw=Ink, line width=0.7pt] (0,0) rectangle (3,3);
  \node[font=\bfseries\small] at (1.5,-0.35) {Reserve A: one block};
  \foreach \x/\y in {4.8/0,6.6/0,4.8/1.8,6.6/1.8}{\fill[black!25, draw=Ink, line width=0.7pt] (\x,\y) rectangle ++(1.5,1.5);
    \draw[Ink, line width=0.5pt, densely dotted] (\x+0.18,\y+0.18) rectangle ++(1.14,1.14);}
  \draw[Ink, line width=0.5pt, densely dotted] (0.18,0.18) rectangle (2.82,2.82);
  \node[font=\bfseries\small] at (6.45,-0.35) {Reserve B: four patches};
\end{tikzpicture}
\figcap{Both reserves have the same area. The dotted lines mark the edge zone, which is hotter, windier and weedier.}
\end{center}

\begin{workedex}{Why the patches lose}
Reserve A is one 40\,ha block; Reserve B is four 10\,ha blocks separated by farmland.

\smallskip
In Reserve B, a much bigger share of the bush is edge, where weeds and foxes get in and conditions are harsh. Animals that need a large area, or must cross open farmland to find mates, cannot survive in a 10\,ha patch. So Reserve B supports fewer native species, even with the same total area.
\end{workedex}

\begin{wedo}{Fragmentation}
A new freeway is built through a forest. Suggest two ways to reduce the harm to the forest's animals, and explain how each helps.
\answerbox{5cm}
\end{wedo}

\subsection{Threats that Combine}\label{sub:10.2}

\begin{workedex}{Leadbeater's possum}
Leadbeater's possum, Victoria's animal emblem, nests only in hollows of mountain ash trees more than about 150 years old.

\smallskip
\textbf{Threat 1, habitat loss:} logging has removed most of the old trees with hollows.\\
\textbf{Threat 2, fire:} the 2009 Black Saturday fires burnt nearly half of its remaining habitat.\\
\textbf{Threat 3, climate change:} fires come more often, so regrowing forest burns again before it can grow hollows.

\smallskip
Each threat on its own is serious; together, they push the possum towards extinction.
\end{workedex}

\begin{wedo}{Threats that combine}
The mountain pygmy possum needs deep snow to hibernate under, and bogong moths to eat. Snow cover is shrinking, and drought has cut moth numbers. Explain why these two threats together are worse than either alone.
\answerbox{5cm}
\end{wedo}

\begin{countrybox}{First Nations perspectives --- A long memory of Country}
Written records in Victoria go back only about 190 years. Wurundjeri Woi-wurrung oral histories and knowledge of Country reach back tens of thousands of years, and remember animals and plants that are now gone from an area. When Elders say a species ``should be here'', that is important information for setting goals to restore an ecosystem.
\end{countrybox}

% ─── PART C ───────────────────────────────────────────────────────────────────
\lessonpart{C}{Practice}{Amber}\label{prac:10}

\task Name the main human impact in each case: habitat loss, fragmentation, pollution, climate change, introduced species or over-harvesting.
\begin{ctab}
\begin{tabular}{|SW{6.2cm}|W{3.8cm}|}
  \hline\hd & \thead{Case} & \thead{Main impact} \\ \hline
  & \scaf{A housing estate is built on native grassland} & \scaf{Habitat loss} \\ \hline
  & Fertiliser runoff causes an algal bloom in the Gippsland Lakes & \\ \hline
  & Snow cover in the Victorian Alps has fallen since 1980 & \\ \hline
  & Rabbits eat the grass native grazers need & \\ \hline
  & A freeway splits a forest into two patches & \\ \hline
  & Too many abalone are taken from Port Phillip Bay & \\ \hline
\end{tabular}
\end{ctab}

\needspace{14\baselineskip}
\task Trace the eutrophication of the Gippsland Lakes, one step per part.
\pt{a} Where do the extra nutrients come from?
\answerbox{2.5cm}
\pt{b} What happens to the algae?
\answerbox{2.5cm}
\pt{c} What happens to the water clarity and the water plants?
\answerbox{2.5cm}
\pt{d} What happens to the oxygen and the fish?
\answerbox{2.5cm}

\needspace{10\baselineskip}
\task Choose one Victorian species facing more than one threat. Name it, give two threats, and explain why the combination is worse than either threat alone.
\answerbox{5cm}

% ─── PART D ───────────────────────────────────────────────────────────────────
\lessonpart{D}{Above level}{Violet}\label{ext:10}

\begin{yr8box}
These questions go beyond the Year 7 standard.

\medskip
\begin{keepext}
\exttask A population of eastern barred bandicoots has been cut off in a 15\,ha reserve for 20 years. Explain two risks that small, isolated populations face. Use the terms \emph{genetic diversity} and \emph{inbreeding}.
\answerbox{6cm}
\end{keepext}
\begin{keepext}
\exttask Explain why losing biodiversity can make it harder for the world to deal with other problems, such as climate change or supplies of fresh water.
\answerbox{5cm}
\end{keepext}
\end{yr8box}

% ══════════════════════════════════════════════════════════════════════════════
%  LESSON 11 — CONSERVATION AND ACTION
% ══════════════════════════════════════════════════════════════════════════════
\setcounter{section}{10}
\sectiondiv{Conservation and Action}{In-situ and ex-situ \textbullet{} A Victorian success story \textbullet{} What we can do}\label{sec:11}
\lessoncodes{VC2S8U04 \textbullet{} VC2S8H03 \textbullet{} VC2S8H04}

\vocabhead
\begin{ctab}
\begin{tabular}{|SW{3.4cm}|W{7.0cm}|}
  \hline\hd & \thead{Word} & \thead{Definition (copy from the slide)} \\ \hline
  & \scaf{\textbf{Conservation}} & \scaf{Protecting species and ecosystems so they survive into the future.} \\ \hline
  & \textbf{In-situ} & \\ \hline
  & \textbf{Ex-situ} & \\ \hline
  & \textbf{Captive breeding} & \\ \hline
  & \textbf{Reintroduction} & \\ \hline
  & \textbf{Habitat corridor} & \\ \hline
\end{tabular}
\end{ctab}

% ─── PART A ───────────────────────────────────────────────────────────────────
\lessonpart{A}{Reading and comprehension}{Teal}\label{read:11}

\begin{readingbox}{Reading --- Back from the Brink}
By 1989 the eastern barred bandicoot was gone from the wild in mainland Victoria. Foxes, cats and the loss of its grassland habitat had wiped it out, except for a few animals in captivity. Its story since then is one of Australia's great \textbf{conservation} successes.

Conservation works in two ways. \textbf{In-situ} conservation (``in place'') protects species where they live: national parks, fences that keep out foxes and cats, and \textbf{habitat corridors}, strips of bush that join separate patches so animals can move between them. \textbf{Ex-situ} conservation (``out of place'') protects species away from the wild: zoos, seed banks, and \textbf{captive breeding}, where animals are bred in safety.

The bandicoot recovery used both. Zoos Victoria and Healesville Sanctuary bred bandicoots in captivity. Meanwhile, predator-free areas were made on the mainland and on fox-free islands, ready for their \textbf{reintroduction}: releasing animals back into the wild. In 2021 it became the first in Australia to be moved from ``extinct in the wild'' to ``endangered''.
\end{readingbox}

\rcreset
\rcq{Why had the eastern barred bandicoot disappeared from the wild?}{2.5cm}
\rcq{What is the difference between in-situ and ex-situ conservation?}{2.5cm}
\rcq{What is a habitat corridor, and how does it help animals?}{2.5cm}
\rcq{What had to be ready before bandicoots could be reintroduced?}{2.5cm}
\rcq{What change in the bandicoot's status happened in 2021?}{2cm}

% ─── PART B ───────────────────────────────────────────────────────────────────
\lessonpart{B}{Theory and examples}{Blue}\label{theory:11}

\subsection{In-situ and Ex-situ}\label{sub:11.1}

\begin{ctab}
\begin{tabular}{|L{4.9cm}|L{4.9cm}|}
  \hline\thead{In-situ (in place)} & \thead{Ex-situ (out of place)} \\ \hline
  National parks and reserves & Captive breeding in zoos \\ \hline
  Habitat corridors & Seed banks \\ \hline
  Fox and cat control & Freezing eggs, sperm or tissue \\ \hline
  Laws protecting native vegetation & Growing rare plants in nurseries \\ \hline
\end{tabular}
\end{ctab}

\begin{workedex}{Evaluating a strategy: habitat corridors}
\textbf{How it works:} planting strips of native bush to join isolated patches.\\
\textbf{Evidence:} connected patches hold more bird species than isolated patches of the same size; powerful owls use creek corridors to move through Melbourne.\\
\textbf{Advantages:} cheaper than buying land; lets animals find mates and food.\\
\textbf{Limitations:} weeds, fire and disease can spread along corridors too; narrow corridors are not used by some species.\\
\textbf{Judgement:} effective, especially for animals that move a lot, but best used alongside reserves and predator control.
\end{workedex}

\begin{wedo}{Evaluating a strategy: captive breeding}
Use the same headings (how it works, an advantage, a limitation, a judgement) to evaluate captive breeding.
\answerbox{6cm}
\end{wedo}

\begin{countrybox}{First Nations perspectives --- The oldest conservation system}
First Nations peoples have cared for Country and its biodiversity for at least 65\,000 years. Caring for Country is a living system of land management, not a thing of the past. Indigenous Protected Areas, managed by First Nations communities, now make up the largest share of Australia's protected land. For Wurundjeri Woi-wurrung people, managing Country again is also about self-determination and healing.
\end{countrybox}

% ─── PART C ───────────────────────────────────────────────────────────────────
\lessonpart{C}{Practice}{Amber}\label{prac:11}

\task Rank these strategies from 1 (most effective) to 6 (least effective), in your opinion, and label each in-situ or ex-situ.
\begin{ctab}
\begin{tabular}{|SW{5.6cm}|Z{2.0cm}|Z{1.8cm}|}
  \hline\hd & \thead{Strategy} & \thead{In-situ or ex-situ} & \thead{Rank} \\ \hline
  & \scaf{Create a new national park} & \scaf{In-situ} & \\ \hline
  & Run a captive breeding program & & \\ \hline
  & Plant corridors between bush patches & & \\ \hline
  & Freeze seeds in a seed bank & & \\ \hline
  & Pass a law banning a harmful pesticide & & \\ \hline
  & Teach the community about pest species & & \\ \hline
\end{tabular}
\end{ctab}
Justify your top choice.
\answerbox{4cm}

\needspace{14\baselineskip}
\task The eastern barred bandicoot recovery program bred bandicoots at zoos, then released them into fox-free enclosures and islands.
\pt{a} Is this in-situ, ex-situ or both? Explain.
\answerbox{4cm}
\pt{b} By 2022 there were more than 1500 bandicoots in the wild. What does this tell you about the program? What worries remain?
\answerbox{5cm}

\needspace{14\baselineskip}
\task Design a conservation action your school or community could take for native wildlife in Melbourne's northern suburbs.
\pt{a} What is the action?
\answerbox{2.5cm}
\pt{b} How would it help native species? Use an idea from this unit.
\answerbox{4cm}
\pt{c} Give one limitation of your action.
\answerbox{3cm}

% ─── PART D ───────────────────────────────────────────────────────────────────
\lessonpart{D}{Above level}{Violet}\label{ext:11}

\begin{yr8box}
These questions go beyond the Year 7 standard.

\medskip
\begin{keepext}
\exttask The EPBC Act is Australia's main law for protecting threatened species. A 2020 review found it ``not fit for purpose'', because biodiversity keeps falling. Explain what the law is meant to do, what evidence suggests it is failing, and one improvement you would make.
\answerbox{6cm}
\end{keepext}
\begin{keepext}
\exttask Some scientists propose bringing Tasmanian devils back to mainland Victoria to control foxes and cats. Give one ecological benefit, one concern farmers or communities might raise, and say who you think should decide, and why.
\answerbox{6cm}
\end{keepext}
\end{yr8box}

\clearpage
{\headingfont\bfseries\Large Final reflection}\par\vspace{0.5mm}\noindent\rule{\linewidth}{0.6pt}\par\vspace{2mm}
What is the one idea from this unit that most changed how you think about the living world? Write a paragraph.
\answerbox{11cm}
